\section{Limitations}
\label{sec:limitations}

This section states plainly what the psl-og member is not, in the spirit of the suite-wide rule that a model's status --- replication, emulation, or structural counterfactual --- is part of its result.

\paragraph{No ground truth.} The OBR emulator and the SVAR members of the suite are validated against published figures; the OLG member has nothing comparable to replicate. Its outputs are structural counterfactuals: internally consistent general-equilibrium answers \emph{conditional on} the model's preferences, technology, demographics, and closure rule. Two defensible OLG models can disagree materially on the long-run effect of the same reform --- the OBR's own UK OLG model \citep{OBR2025OLG} provides a natural cross-check, and side-by-side comparison against it is the closest available substitute for validation --- so point estimates should be read as model-consistent magnitudes and signs, not forecasts.

\paragraph{Experimental-grade wiring.} The integration is young. The version-pin fragility of Section~\ref{subsec:pins} means the OG member currently runs against \texttt{policyengine-uk} 2.88.0 while the suite's microsimulation member tracks the latest release, so the two members can disagree about baseline statute; the gated enhanced-FRS microdata adds an access prerequisite; and the CI cannot exercise a seventeen-minute solve on every commit, so regressions in solve behaviour are caught by scheduled runs and by hand, not continuously.

\paragraph{Steady-state-only scoring today.} The CLI compares long-run steady states. That is not a budget-window costing: it answers ``where does the economy settle?'', not ``what happens in each of the next five fiscal years?''. The transition-path solver exists and is reachable through the \texttt{oguk} API, but it is hours-long, Dask-distributed, and not yet wired into \texttt{pe-macro score}. Every scoring payload carries this caveat explicitly in its \texttt{ScoreResult} block.

\paragraph{Fast-configuration compromises.} The deployed configuration pools tax functions across ages and uses one production sector. Pooling smooths over pension-age discontinuities in effective schedules; a single sector rules out compositional questions. Both dials can be turned up (\texttt{brackets}/\texttt{each}, \texttt{multi\_sector=True}) at material runtime and stability cost --- OLG solves can fail to converge on extreme reforms, and the documented mitigation (raise \texttt{max\_iter}, fall back to pooled) is a workaround, not a guarantee.

\paragraph{Calibration vintage and inheritance.} Calibration inputs carry vintages: demographic profiles from a UN release, depreciation and factor shares from a Blue Book vintage, potential growth from a particular EFO, and --- most importantly --- lifetime-ability profiles $e_{j,s}$ inherited from the US-estimated OG-Core defaults rather than re-estimated on UK microdata. No published achieved-versus-target moment reconciliation exists for the deployed version. Results that lean on the upper tail of the earnings distribution, or on fine age detail around retirement, are the least trustworthy outputs of the current calibration.

\paragraph{Model-class assumptions.} Finally, the standard structural caveats apply: perfect foresight (no aggregate risk), rational optimising households, no involuntary unemployment, exogenous productivity growth, and a closure rule that mechanically stabilises debt. These are the assumptions that buy general-equilibrium discipline; conclusions that would reverse under plausible relaxations of them should not be attributed to the data.
